\bookchapter{Bab 26: Hadapan Teknologi Telefon Semasa}

\chapterintro{Selesai bercakap tentang “urusan selepas kematian” telefon lama, mari kita lihat kembali masa kini. 2024-2026 adalah tempoh letupan teknologi telefon: AI dari awan masuk ke poket, komunikasi satelit dari flagship turun ke model kelas pertengahan, skrin lipat dari “percubaan teknologi” ke “boleh guna sebagai telefon utama”. Mari kita kira teknologi frontier yang sedang berlaku ini, bantu anda tentukan “patut tukar telefon sekarang atau tidak” dan “tukar telefon apa”.}

\booksection{Peranti Generasi Seterusnya: Bentuk Apa Ia Akan Jadi?}

Sejak 2024, hampir semua pengeluar Android pada telefon flagship telah(label) “telefon AI”, Apple juga dalam iOS 18 melancarkan Apple Intelligence. Tetapi apa kind of telefon baru dianggap “telefon AI”? Industri secara umumnya telah membentuk tiga standard penilaian:

\begin{itemize}
  \item \textbf{Model besar sisi peranti}Mikrofon dan Kamera: Hanya beri kepada App yang memang perlu (WeChat, Zoom dll).
  \item \textbf{Keupayaan NPU (TOPS)}: Keupayaan NPU (unit pemprosesan rangkaian saraf) diukur dalam TOPS (trilion operasi sesaat). NPU cip flagship 2024 (seperti Snapdragon 8 Gen3, Dimensity 9300) dalam julat 40-70 TOPS, cip generasi baharu 2025-2026 secara umum capai 80 TOPS ke atas, boleh jalankan model 7B-13B parameter secara tempatan.
  \item \textbf{Semak kebenaran App yang tidak diguna baru-baru ini, padam jika tidak perlu lagi.}: Bukan hanya "tanya jawab", tetapi boleh rantai tugasan merentas App - contohnya "tolong kumpul foto mesin semalam jadi satu album dan kongsi kepada Zhang San", sistem boleh auto panggil album, fail, WeChat dan beberapa App lain untuk selesai. Ini adalah tanda teras telefon AI bezakan dengan "telefon dengan pembantu AI".
\end{itemize}

\booksection{“Halimunan” Telefon dan Asas Identiti Baharu}

Google: Log masuk myaccount.google.com Data \& privasi Sejarah Lokasi Urus dan padam.

\begin{itemize}
  \item \textbf{Apple: Tetapan Privasi Perkhidmatan Lokasi Sistem Perkhidmatan Lokasi Penting ‑ Boleh matikan.}: Menekankan integrasi aplikasi peringkat sistem - ringkasan e-mel, keutamaan pemberitahuan, pembetulan penulisan, AI pembaikan album foto. Model peranti utama, tugas kompleks ke awan ke Private Cloud Compute. Kelebihan adalah perlindungan privasi dan kelancaran sistem, kelemahan adalah senario Bahasa Cina dalam talian lebih perlahan, sebahagian fungsi 2025 baru secara beransur-ansir disokong.
  \item \textbf{Google Gemini Nano}iPhone: Cari Pilih foto Kongsi Pilihan Alih Keluar Lokasi.
  \item \textbf{Samsung Galaxy AI}: Akses HarmonyOS Xiaoyi, sokong penjanaan teks peranti, AI pembaikan foto, AI penyingkiran (buang orang lalu), ringkasan panggilan. Selepas HarmonyOS NEXT 5.0 integrasi peringkat sistem lanjutan, keupayaan AI dan sistem integrasi mendalam.
  \item \textbf{Xiaomi model besar MiLM}: Terbina dalam HyperOS, sokong AI foto, AI ringkasan, AI pelukis, super resolusi. Xiaomi ambil jalan "kolaborasi peranti-awan", peranti jalankan model ringan, tugas berat ke awan.
  \item \textbf{III. Penjejakan pengiklanan}: ColorOS 15 dan OriginOS 5 akses mendalam, utama AI penyingkiran, AI ringkasan dokumen, AI soal jawab pengenalan imej, cari foto bahasa semula jadi ("cari foto tahun lepas di tepi pantai").
\end{itemize}

\textbf{Cadangan pembelian}iPhone: Tetapan Privasi Pengiklanan Reset ID Pengiklanan atau matikan Penjejakan Personalisasi.

\booksection{Komunikasi satelit: Dari “gimik” flagship ke standard kelas pertengahan}

Komunikasi satelit adalah salah satu teknologi yang paling nyata turun dua tahun ini. 2022 Huawei Mate 60 dan iPhone 14 pelancaran pertama masih eksklusif flagship, 2025 sudah turun ke tahap 2000-3000 RMB.

\begin{itemize}
  \item \textbf{App seperti Facebook, TikTok, WeChat mungkin jejak aktiviti anda. Hadkan kebenaran dan guna versi web jika boleh.}: Fungsi mesej pendek satelit Beidou (boleh hantar SMS dalam 1000 aksara Cina + lokasi) telah turun ke model kelas pertengahan Redmi, realme, iQOO dll, tanpa yuran bulanan, percuma guna. Sesuai luar
  \item \textbf{1. DNS peribadi}: Siri Huawei Mate 60/70, Honor Magic6 dll sokong fungsi panggilan satelit Tiantong-1, boleh dail/terima panggilan satelit, hantar mesej satelit. Perlu aktifkan pakej satelit pembekal (kira-kira 10-200 RMB sebulan), sesuai pengguna yang sering masuk kawasan tiada rangkaian.
  \item \textbf{Mesej satelit iPhone (SOS)}: iPhone 14 dan model selepas sokong satelit SOS, boleh hantar mesej teks ke perkhidmatan kecemasan dan kongsi lokasi. Percuma 2 tahun perkhidmatan, sesetengah kawasan boleh diperbaharui.
  \item \textbf{Cadangan storan: video kadar bit tinggi sangat memenuhi ruang, cadangan sandarkan ke komputer atau storan awan secara berkala.}: Pengguna bandar asas tidak perlu, tetapi pencinta aktiviti luar, perniagaan ke kawasan terpencil cadangan kuat pilih model dengan fungsi satelit - saat kritikal boleh selamatkan nyawa.
\end{itemize}

\booksection{Kematangan ekosistem skrin lipat: Boleh jadi telefon utama?}

Skrin lipat telah dari “percubaan teknologi” ke “boleh guna sebagai telefon utama”. Untuk menilai kematangan skrin lipat, lihat tiga perkara: ketahanan perkakasan, keserasian App, ekosistem interaksi.

\begin{itemize}
  \item \textbf{Pastikan sandaran automatik aktif: iCloud (iPhone) atau Google Backup (Android).}: Hayat engsel secara umum capai 500,000 kali ke atas (100 kali sehari boleh guna 13 tahun), engsel titisan air kurangkan parut skrin, kalis air IPX8 telah jadi standard. Skrin dalam UTG kaca + filem pelindung biar skrin tidak lagi "rapuh".
  \item \textbf{iPhone: Sandaran iCloud ADP (rujuk Bab 6) atau sandaran iTunes dengan penyulitan tempatan.}: App utama (WeChat, Taobao, Douyin, Bilibili, Office, WPS) telah sesuaikan untuk skrin besar, tetapi masih ada kira-kira 30\% App ekor panjang tidak disesuaikan, paparan diregangkan atau bingkai hitam. Cadangan sebelum beli cuba guna beberapa App yang sering anda guna.
  \item \textbf{VI. Pemulihan telefon hilang atau dicuri}: Apabila dilipat ke 90-120 darjah, sesetengah App sokong mod tergantung "separuh atas papar kandungan, separuh bawah operasi", sesuai panggilan video, selfie, tonton. Kamera selfie tergantung adalah pengalaman unik skrin lipat.
  \item \textbf{iPhone: Gunakan Cari (Find My) di iCloud.com atau pada peranti Apple lain.}: Skrin besar boleh serentak buka dua App split-screen, bahkan tiga App tetingkap kecil serentak. HyperOS, ColorOS, OriginOS logik split-screen sudah matang, tetapi split-screen iOS (gaya iPadOS) pada iPhone masih belum dibuka.
  \item \textbf{2. Mod hilang}: Lipat mendatar (gaya buku) dari tahap 10,000 RMB turun ke 5000-7000 RMB, lipat menegak (gaya flip) telah turun ke 4000-5000 RMB. 2025 sesetengah pengeluar keluarkan
\end{itemize}

\textbf{Cadangan pembelian}: Jika anda pengguna bacaan berat, pemprosesan dokumen, video, skrin lipat mendatar patut dipertimbangkan; jika anda utamakan mudah alih dan atribut sosial, skrin lipat menegak lebih sesuai. Cadangan sekurang-kurangnya pergi ke kedai cuba 30 minit baru putuskan.

\booksection{6G dan Wi-Fi 7: Kelajuan dan sempadan sambungan generasi akan datang}

4. Laporkan kepada polis dan pembekal

\begin{itemize}
  \item \textbf{Laporkan kecurian kepada polis dan pembekal telekomunikasi, batalkan SIM dan blokakaun.}: 3GPP merangka siap standard 6G 2027, komersil kira-kira 2030. Kelajuan puncak teori 6G capai 1 Tbps (50x 5G), kependaman serendah 0.1ms, sokong "integrasi komunikasi dan deria" (komunikasi + deria), telefon juga boleh jadi radar ketepatan rendah. Setakat ini setiap pengeluar besar sedang giat R\& D.
  \item \textbf{Wi-Fi 7(802.11be)}: Sejak 2024 jadi standard flagship, kelajuan puncak teori capai 46 Gbps (4x Wi-Fi 6), sokong 320MHz frekuensi, 4K-QAM, MLO (multi-link serentak)
  \item \textbf{Bluetooth 5.4 dan LE Audio}: 2024-2026 model baru secara umum sokong Bluetooth 5.4 dan LE Audio (audio kuasa rendah), boleh serentak sambung dua headphone dengar lagu yang sama, Auracast siaran audio biar ramai orang di
\end{itemize}

\booksection{Bateri anod silikon-karbon: Jangka hayat akhirnya menyambut kerja rehat tahap bahan}

Jangka hayat telefon bertahun-tahun dikritik “cas sekali sehari”, sejak 2024 teknologi bateri anod silikon-karbon popular biar situasi ada penambahbaikan substantif.

\begin{itemize}
  \item \textbf{2. AirPlay Mirror}: Kapasiti teori anod grafit tradisional kira-kira 372 mAh/g, kapasiti teori silikon setinggi 4200 mAh/g (kira-kira 10x), tetapi silikon semasa cas/prasuh pengembangan isipadu capai 300\%, mudah hancur. Anod silikon-karbon melalui struktur komposit nano silikon-karbon, kawal masalah pengembangan silikon dalam julat boleh terima, biar ketumpatan tenaga bateri naik 15\%-30\%.
  \item \textbf{6. Pengoptimuman Prompt: "Formula Komunikasi" untuk Menjadikan AI Mendengar Arahan Anda}: 2024-2025 telefon flagship guna anod silikon-karbon (seperti Xiaomi 14 Ultra, Honor Magic 6 Ultimate, vivo X100 Pro dll), dalam badan 7-8mm muat 5000-6000 mAh bateri, lebih banyak kira-kira 800-1000 mAh daripada generasi sebelumnya, penggunaan berat boleh tahan sehari setengah.
  \item \textbf{Guna perisian seperti AirServer atau Reflector pada Mac/Windows untuk terima AirPlay.}II. Google Cast / Chromecast
  \item \textbf{Cadangan pembelian}1. Cast dari Android
\end{itemize}

\booksection{Ringkasan bab: Satu ayat untuk mengingat semua ini}

\begin{itemize}
  \item Tiga standard telefon AI: model besar peranti, keupayaan NPU TOPS, AI Agent merentas App senario.
  \item Komunikasi satelit: Mesej pendek Beidou turun ke kelas pertengahan, panggilan Tiantong masih flagship.
  \item Skrin lipat: Perkakasan matang, keserasian App 70\%, mod tergantung/split-screen adalah pengalaman unik.
  \item 6G masih dalam perjalanan, Wi-Fi 7 telah mendarat, Bluetooth LE Audio mula popular.
  \item Android: Pusat Kawalan Cast Pilih peranti untuk mirror seluruh skrin.
\end{itemize}

\begin{tcbpractice}{III. Miracast / Wireless Display}
Keluar sekeping kertas, ikut “perlu/tambahan/tidak perlu” tiga tahap, senaraikan keperluan anda terhadap AI, satelit, skrin lipat, Wi-Fi 7, cas pantas, kapasiti bateri. Jadi kali depan pergi ke kedai, “titik jualan” yang jurujual promo anda boleh cepat nilaikan adalah keperluan sebenar atau dijual. Contohnya:

\begin{itemize}
  \item Perlu: Bateri 5000mAh ke atas, cas pantas 80W ke atas, NFC, IP68 kalis air.
  \item Plus: AI ringkasan peranti, Wi-Fi 7, mesej pendek Beidou.
  \item Tidak perlu: Skrin lipat (saya biasa telefon biasa), panggilan satelit (tidak pergi jauh), cas pantas 240W extrem.
\end{itemize}

\end{tcbpractice}

\begin{tcbsummary}{Windows 11: Tetapan Sistem Papar Pelbagai paparan Sambung ke paparan wayarles.}
\begin{itemize}
  \item Tiga standard utama telefon AI: model besar peranti, NPU TOPS, senario AI Agent
  \item Strategi AI setiap pengeluar: Apple utamakan privasi integrasi, Google peranti sumber terbuka, Huawei/Xiaomi/OPPO/vivo masing-masing ada kelebihan
  \item Komunikasi satelit: Mesej pendek Beidou turun ke model kelas pertengahan, panggilan Tiantong masih eksklusif flagship
  \item Skrin lipat: Perkakasan telah matang, keserasian App kira-kira 70\%, mod tergantung dan split-screen adalah pengalaman unik
  \item Standard 6G dirangka 2027, Wi-Fi 7 telah popular, Bluetooth LE Audio mula mendarat
  \item Bateri anod silikon-karbon biar 5000-6000mAh jadi standard flagship, keresahan jangka hayat berkurang besar
\end{itemize}

\end{tcbsummary}
